\documentclass[10pt,a4paper]{amsart}
\usepackage[margin=19mm]{geometry}
\usepackage{amsmath,amssymb,mathtools}
\usepackage[scr=rsfs,cal=euler]{mathalfa}
\usepackage{xfrac}
\usepackage[unicode,hidelinks,bookmarks=false]{hyperref}
\newcommand{\ie}{i.e.\ }
\newcommand{\cf}{cf.\ }
\newcommand{\resp}{respectively\ }
\newcommand{\Subsection}{\subsection}
\providecommand{\nobreakdash}{\nobreak\hbox{-}}
\setlength{\emergencystretch}{2em}
\multlinegap0pt
\usepackage{amssymb}
\usepackage{xcolor}
\usepackage{algorithm}\usepackage{algpseudocode}
\usepackage{tikz}
\usepackage{mathtools}
\newtheorem{lemma}{Lemma}
\title{Regularity and singularity of the blow-up curve for a wave equation with a derivative nonlinearity and a scale-invariant damping}
\author{Ahmed Bchatnia, Makram Hamouda, Firas Kaabi, Takiko Sasaki, Hatem Zaag}
\date{Source: arXiv:2604.03848v1 (CC BY 4.0)}
\begin{document}
\maketitle

\noindent\emph{Source and licence.} Regularity and singularity of the blow-up curve for a wave equation with a derivative nonlinearity and a scale-invariant damping. Version arXiv:2604.03848v1 (CC BY 4.0), distributed by arXiv under the Creative Commons Attribution 4.0 International licence, http://creativecommons.org/licenses/by/4.0/, as recorded in the arXiv licence metadata for this exact version and shown on the abstract page https://arxiv.org/abs/2604.03848v1. Full licence notice: \texttt{licenses/arxiv-2604.03848-CC-BY-4.0.txt}.\par\smallskip

\noindent\textit{Editorial scope.} Counted unit: the article's lemma 2.3 (source0.tex lines 597--599). The article's complete original proof of it is delivered with it (source0.tex lines 602--809, one \texttt{proof} environment of 208 lines). No mathematics is written, repaired or re-derived: the statement and the proof are the article's own lines, in the article's own notation, and no retained line is substituted or re-flowed.  Retained uncounted as the source-local context the proof needs: lines 246--253; lines 477--483; lines 573--591.  Nothing was omitted from any retained span, so the package carries no omission record.  No retained line contains a citation, so the wrapper carries no bibliography and no bibliography entry.  No retained line contains a figure environment, an image-inclusion command or drawing code, so no image is delivered and the package declares no asset dependency.  Editorial formatting only, in addition: an AMS \texttt{amsart} preamble that restates the article's own declarations and macros used by the retained lines, namely the environments \texttt{lemma} on separate counters, together with the standard packages those lines need.  The retained environments print in this document as Lemma 1 in order of appearance, whereas the article prints the same items with the numbering of its own full document.  The complete native source of the article as distributed in the arXiv e-print for this exact version is bundled unchanged as \texttt{sources/197840/source0.tex}.
\begin{equation}
\label{1.5}
\begin{cases}
D_- \phi = 2^{-p}|\phi + \psi|^p - \dfrac{\mu}{1 + t} \dfrac{\phi + \psi}{2}, & x \in \mathbb{R}, t > 0, \\
D_+ \psi = 2^{-p}|\phi + \psi|^p - \dfrac{\mu}{1 + t} \dfrac{\phi + \psi}{2}, & x \in \mathbb{R}, t > 0, \\
\phi(x, 0) = f(x), \quad \psi(x, 0) = g(x), & x \in \mathbb{R},
\end{cases}
\end{equation}

\begin{lemma}[Monotonicity of Approximations]\label{lem:2.1}
Under assumption $(A_2)$, the successive approximations satisfy:
$$
\phi_{n+1}(x,t) \geq \phi_n(x,t) \geq 0 \quad \text{and} \quad \psi_{n+1}(x,t) \geq \psi_n(x,t) \geq 0,
$$
for all $(x,t) \in K_{R^*,T^*}$ and every $n \in \mathbb{N} \cup \{0\}$. Moreover, this monotonicity holds uniformly across the domain.
\end{lemma}

Let $(x,t) \in K_{R^*, T^*}$. Since $\{\phi_n(x,t)\}_{n=0}^\infty$ and $\{\psi_n(x,t)\}_{n=0}^\infty$ are increasing sequences in $n$, we have
\begin{equation}
\left\{
  \begin{aligned}
    \lim_{n \to \infty} \phi_n(x,t) &= \sup_{n \in \mathbb{N}} \phi_n(x,t) = \phi(x,t), \\
    \lim_{n \to \infty} \psi_n(x,t) &= \sup_{n \in \mathbb{N}} \psi_n(x,t) = \psi(x,t).
  \end{aligned}
    \label{2.3}
    \right.
\end{equation}
From Lemma~\ref{lem:2.1}, it follows that $\phi$ and $\psi$ increase monotonically in $t$. Hence, there exists a function $T(x)$ such that
$$
  T(x) = \sup \{\, t \in (0, T^*) : (\phi + \psi)(x,t) < \infty \,\}, \qquad x \in B_{R^*}.
$$
We define
\begin{equation}
    \label{omega}
  \Omega = \{\, (x,t) : x \in B_{R^*}, \; 0 < t < T(x) \,\}.
\end{equation}

\begin{lemma}\label{lemma 2.3}
  Under the assumptions $(A_2)$--$(A_4)$, the pair $(\phi, \psi)$,  defined by \eqref{2.3}, is the unique solution of \eqref{1.5} in the class $(\mathcal{C}^{3,1}(\Omega))^2$.
\end{lemma}

\begin{proof}
  Define the set
  $$
    B(t) = \bigl\{ x \in B_{R^*+T^*}: \lvert x - \tilde{x} \rvert \le \tilde{t}-t \bigr\}, \quad (\tilde{x}, \tilde{t}) \in \Omega.
  $$

  \medskip
  \noindent\textbf{(Proof of regularity).}  
  We first show that \((\phi, \psi)\) is a \(\mathcal{C}^{3,1}(\overline{\Omega})^2\) solution of \eqref{1.5}.  
  The proof is carried out in two steps.

  \smallskip
  \noindent\textbf{Step 1.}  
  Fix \((x, \tilde{t}\,) \in \Omega\).  
  We claim that there exists a constant \(M_0 > 0\) such that
  \begin{equation}
    \label{eq:2.4}
    \lVert \phi + \psi \rVert_{L^\infty(B(t))} \le M_0, \quad \text{for all } t \in [0, \tilde{t}\,].
  \end{equation}

  \noindent We argue by contradiction. Define  
  $$
    Y_{\tilde{x}} = \{ x \in B_{R^*} : \lvert x - \tilde{x} \rvert \le \tilde{t} - T(x) \},
  $$
  where \(m\) denotes the 1-dimensional Lebesgue measure.  
  Suppose \eqref{eq:2.4} fails, i.e., \(m(Y_{\tilde{x}}) > 0\).  
  By the monotonicity of \(\phi_n + \psi_n\) in \(t\), we have  
  $$
    (x, t) \notin \Omega \quad \text{if } x \in Y_{\tilde{x}} \text{ and } t = x + \tilde{t} - \tilde{x} \text{ or } t = -x + \tilde{t} + \tilde{x}.
  $$
  Hence, either \(m(Y_{\tilde{x},+}) > 0\) or \(m(Y_{\tilde{x},-}) > 0\), where  
  $$
    Y_{\tilde{x},+} = \{ s \in (0, \tilde{t}) : s = x + \tilde{t} - \tilde{x}, \ x \in Y_{\tilde{x}} \},  
  $$
  $$
    Y_{\tilde{x},-} = \{ s \in (0, \tilde{t}) : s = -x + \tilde{t} + \tilde{x}, \ x \in Y_{\tilde{x}} \}.
  $$
  However, evaluating the iteration yields  
  $$
    \infty > (\phi_{n+1} + \psi_{n+1})(\tilde{x}, \tilde{t}) 
    \geq \int_{Y_{\tilde{x},+}} \left( 2^{-p} \lvert \phi_n + \psi_n \rvert^p - \frac{\mu}{1+t} \frac{\phi_n + \psi_n}{2} \right)(x + \tilde{t} - s, s) \, ds  
  $$
  $$
    + 2^{-p} \int_{Y_{\tilde{x},-}} \left( 2^{-p} \lvert \phi_n + \psi_n \rvert^p - \frac{\mu}{1+t} \frac{\phi_n + \psi_n}{2} \right)(-x + \tilde{t} + \tilde{x}, s) \, ds \to \infty \quad \text{as } n \to \infty,
  $$
  which is a contradiction. Thus, \eqref{eq:2.4} holds.

   \noindent\textbf{Step 2.} 
We now prove that $(\phi, \psi)$ is a $\mathcal{C}^{3,1}(\overline{\Omega})^2$ solution of \eqref{1.5}.  

Fix $(\tilde{x}, \tilde{t}) \in \Omega$ and set $K = K(x, \tilde{t})$.  
By localizing the problem, it suffices to show that $(\phi, \psi)$ is a $\mathcal{C}^{3,1}(\overline{K})^2$ solution of \eqref{1.5}.  

First, we establish the uniform convergence of $\phi_n$ to $\phi$ and $\psi_n$ to $\psi$ in $K = K(x, \tilde{t})$.  

By Step~1, there exists a constant $C_0 = C_0(x, \tilde{t}) > 0$ such that
\begin{equation}
  \label{eq:2.51}
  \|\phi_n + \psi_n\|_{L^\infty(B(t))} \leq C_0 \quad \text{for all } t \in [0, \tilde{t}\,] \text{ and } n \in \mathbb{N}.
\end{equation}

For $t \in [0, \tilde{t}\,]$ and $n \in \mathbb{N}$, we estimate the differences:
\begin{align*}
&\|\phi_{n+1} - \phi_n\|_{L^\infty(B(t))} + \|\psi_{n+1} - \psi_n\|_{L^\infty(B(t))} \\
&\quad \leq (p C_0^{p-1} + \mu) \int_0^t \left( 
\|\phi_n - \phi_{n-1}\|_{L^\infty(B(s_1))} + \|\psi_n - \psi_{n-1}\|_{L^\infty(B(s_1))} 
\right) ds_1 \\
&\quad \leq (p C_0^{p-1} + \mu)^2 \int_0^t \int_0^{s_1} \left( 
\|\phi_{n-1} - \phi_{n-2}\|_{L^\infty(B(s_2))} + \|\psi_{n-1} - \psi_{n-2}\|_{L^\infty(B(s_2))} 
\right) ds_2 ds_1.
\end{align*}

Iterating this estimate yields
$$
\|\phi_{n+1} - \phi_n\|_{L^\infty(B(t))} + \|\psi_{n+1} - \psi_n\|_{L^\infty(B(t))} 
\leq 4 C_0 \frac{\big((p C_0^{p-1} + \mu) T\big)^n}{n!} \to 0 \quad \text{as } n \to \infty.
$$

From \eqref{2.3}, we conclude the uniform convergence:
$$
\|\phi_n - \phi\|_{L^\infty(K_{x, \tilde{t}})} + \|\psi_n - \psi\|_{L^\infty(K_{x, \tilde{t}})} \to 0 \quad \text{as } n \to \infty.
$$


We now establish that $\phi, \psi \in W^{1,\infty}(K, (\tilde{x}, \tilde{t}))$. 

The directional derivatives satisfy the following recursive relations for $n \in \mathbb{N} \cup \{0\}$:
$$
\begin{cases}
D_- D_\theta \phi_{n+1} = p 2^{-p} (\phi_n + \psi_n)^{p-1} (D_\theta \phi_n + D_\theta \psi_n), \\
D_+ D_\theta \psi_{n+1} = p 2^{-p} (\phi_n + \psi_n)^{p-1} (D_\theta \phi_n + D_\theta \psi_n),
\end{cases}
$$
where $D_\theta v = (\cos \theta \, \partial_x + \sin \theta \, \partial_t)v$ is the directional derivative.

The initial conditions are given by:
$$
\begin{cases}
D_\theta \phi_{n+1}(x,0) = 
\begin{cases}
(\cos\theta+\sin\theta)\,f^\prime(x) + \sin\theta\cdot 2^{-p}(\gamma_1+\gamma_2)^p & (n = 0), \\
(\cos\theta+\sin\theta)\,f^\prime(x) + \sin\theta\cdot 2^{-p}(f+g)^p(x) & (n \geq 1),
\end{cases} \\[6pt]
D_\theta \psi_{n+1}(x,0) = 
\begin{cases}
(\cos\theta-\sin\theta)\,g^\prime(x) + \sin\theta\cdot 2^{-p}(\gamma_1+\gamma_2)^p & (n = 0), \\
(\cos\theta-\sin\theta)\,g^\prime(x) + \sin\theta\cdot 2^{-p}(f+g)^p(x) & (n \geq 1).
\end{cases}
\end{cases}
$$
\noindent Define the function
$$
W(t) = (C_0^p + \mu C_0) \exp\left((pC_0^{p-1} + \mu + 1) t\right),
$$
which satisfies the integral equation
\begin{equation}
W(t) = C_0^p + \mu C_0 + \int_0^t (pC_0^{p-1} + \mu + 1) W(s) \, ds.
\label{eq:2.4a}
\end{equation}

We prove by induction that for all $t \in [0, \tilde{t}\,]$ and $n \in \mathbb{N} \cup \{0\}$:
\begin{equation}
\|D_\theta \phi_n(\cdot, t)\|_{L^\infty(B(t))} < W(t), \quad 
\|D_\theta \psi_n(\cdot, t)\|_{L^\infty(B(t))} < W(t).
\label{eq:2.5}
\end{equation}

First, observe that $D_\theta \phi_0 = D_\theta \psi_0 = 0 \leq W(t)$ for all $t \geq 0$. Assuming the inequalities hold for $n$, we establish the bound for $n+1$, that is
\begin{equation}
\|2^{-p} (\phi_n + \psi_n)^{p-1} (\cdot, t) (D_\theta \phi_n + D_\theta \psi_n) (\cdot, t)\|_{L^\infty(B(t))} \leq C_0^{p-1} W(t).
\label{eq:2.6}
\end{equation}
From assumption $(A_4)$, we derive the key estimate:
\begin{equation}
\begin{aligned}
\|D_\theta \phi_{n+1}(\cdot, t)\|_{L^\infty(B(t))} 
&\leq 2 \|f^\prime\|_{L^\infty(B(0))} + 2^{-p} \|f + g\|_{L^\infty(B(0))}^p + \frac{\mu}{2} \|f + g\|_{L^\infty(B(0))} \\
&\quad + \int_0^t p \left\| 2^{-p} (\phi_n + \psi_n)^{p-1} (D_\theta \phi_n + D_\theta \psi_n) \right\|_{L^\infty(B(s))} ds \\
&\quad + \int_0^t \left\| \frac{-\mu}{2(1+s)} (D_\theta \phi_n + D_\theta \psi_n) \right\|_{L^\infty(B(s))} ds \\
&\quad + \int_0^t \left\| \frac{\sin \theta}{2} \cdot \frac{\mu}{(1+s)^2} (\phi_n + \psi_n) \right\|_{L^\infty(B(s))} ds \\
&\leq C_0^p + \mu C_0 + \int_0^t (p C_0^{p-1} + \mu + 1) W(s) \, ds = W(t),
\end{aligned}
\label{eq:2.7}
\end{equation}
for $t \in [0, \tilde{t}\,]$. The same bound holds for $\|D_\theta \psi_{n+1}(\cdot, t)\|_{L^\infty(B(t))}$.

\medskip

Define $C_1 = (C_0^p+\mu C_0) \exp((pC_0^{p-1}+\mu+1) T)$, yielding:
\begin{equation}
\|D_\theta \phi_n(\cdot, t)\|_{L^\infty(B(t))} \leq C_1, \quad \text{and} \quad \|D_\theta \psi_n(\cdot, t)\|_{L^\infty(B(t))} \leq C_1.
\label{eq:2.8}
\end{equation}
The estimate of the difference estimates shows that
\begin{equation}
\begin{aligned}
&\|D_\theta \phi_{n+1} - D_\theta \phi_n\|_{L^\infty(B(t))} + \|D_\theta \psi_{n+1} - D_\theta \psi_n\|_{L^\infty(B(t))} \\
&\leq \int_0^t (\mu + p C_0^{p-1}) \big(\|D_\theta \phi_n - D_\theta \phi_{n-1}\|_{L^\infty(B(s_1))} + \|D_\theta \psi_n - D_\theta \psi_{n-1}\|_{L^\infty(B(s_1))}\big) ds_1 \\
&\quad + \int_0^t \big(\mu + 2p(p - 1) C_1 C_0^{p-2}\big) \big(\|\phi_n - \phi_{n-1}\|_{L^\infty(B(s_1))} + \|\psi_n - \psi_{n-1}\|_{L^\infty(B(s_1))}\big) ds_1 \\
&\leq 4C_1 \frac{\big((\mu + p C_0^{p-1}) T\big)^n}{n!} + 4 C_0 C_2 T^n \big(p C_0^{p-1}\big)^{n-1} \sum_{j=1}^n \frac{1}{j! (n - j)!} \to 0,
\end{aligned}
\label{eq:2.91}
\end{equation}
where $C_2 = \mu+2p(p - 1)C_1C_0^{p-2}$.

\medskip

The uniform convergence implies the existence of limits:
$$
\big(\phi_\theta^{(1)}, \psi_\theta^{(1)}\big) \in \big(L^\infty(K_-(\tilde{x}, \tilde{t}))\big)^2,
$$
with
$$
\|D_\theta \phi_n - \phi_\theta^{(1)}\|_{L^\infty(K_-(\tilde{x}, \tilde{t}))} + \|D_\theta \psi_n - \psi_\theta^{(1)}\|_{L^\infty(K_-(\tilde{x}, \tilde{t}))} \to 0.
$$

This establishes $\big(\phi, \psi\big) \in \big(W^{1, \infty}(K_-(\tilde{x}, \tilde{t}))\big)^2$. Iterating the argument yields the higher regularity:
$$
\big(\phi, \psi\big) \in \big(W^{4, \infty}(K_-(\tilde{x}, \tilde{t}))\big)^2 = \big(\mathcal{C}^{3, 1}(K_-(\tilde{x}, \tilde{t}))\big)^2.
$$


\noindent\textbf{Proof of Uniqueness.}
We now establish the uniqueness of the solution $(\phi, \psi)$ to \eqref{1.5}. 
 
Suppose there exist two solutions $(\phi_1, \psi_1)$ and $(\phi_2, \psi_2)$ with corresponding blow-up curves $T_1$ and $T_2$. Define their domains:
$$
\Omega_j = \{(x, t) : x \in B_{R^*}, \ 0 < t < T_j(x)\} \quad \text{for } j = 1,2.
$$
For any $(\tilde{x}, \tilde{t}) \in \Omega_1 \cap \Omega_2$, we have $K_-(\tilde{x}, \tilde{t}) \subset \Omega_1 \cap \Omega_2$. Following the approach in Step 2, we obtain the estimate:
$$\begin{aligned}
&\sup_{0 \leq t' \leq \tilde{t}} \Big( \|\phi_1 - \phi_2\|_{L^\infty(B(t'))} + \|\psi_1 - \psi_2\|_{L^\infty(B(t'))} \Big)\\
&\leq (t p C_0^{p-1} + \mu) \sup_{0 \leq t' \leq \tilde{t}} \Big( \|\phi_1 - \phi_2\|_{L^\infty(B(t'))} + \|\psi_1 - \psi_2\|_{L^\infty(B(t'))} \Big).
\end{aligned}
$$
For sufficiently small $t > 0$, this implies:
$$
\sup_{0 \leq t' \leq t} \Big( \|\phi_1 - \phi_2\|_{L^\infty(B(t'))} + \|\psi_1 - \psi_2\|_{L^\infty(B(t'))} \Big) = 0.
$$
Since $C_0$ is time-independent, we can iterate this argument to extend the uniqueness to all $t' \in [0, \tilde{t}\,]$:
$$
\sup_{0 \leq t' \leq \tilde{t}} \Big( \|\phi_1 - \phi_2\|_{L^\infty(B(t'))} + \|\psi_1 - \psi_2\|_{L^\infty(B(t'))} \Big) = 0.
$$
This establishes that
 $(\phi_1, \psi_1) \equiv (\phi_2, \psi_2)$ in $\Omega_1 \cap \Omega_2$ and
 $T_1(x) = T_2(x)$ for all $x \in B_{R^*}$,
completing the proof of uniqueness and hence the lemma.
\end{proof}

\end{document}
